\documentclass[12pt]{article}

% Packages
\usepackage{geometry}
\geometry{letterpaper, margin=1in}
\usepackage{amsmath}
\usepackage{array}
\usepackage{titlesec}
\usepackage{booktabs}
\usepackage{hyperref}
\usepackage{caption}
\usepackage{fancyhdr}

% Header and Footer
\pagestyle{fancy}
\fancyhf{}
\fancyhead[L]{\textit{University of South Carolina}}
\fancyhead[R]{\thepage}
\fancyfoot[C]{}

% Title Information
\title{Looking to Nature: Nature-Inspired Solutions}
\author{
    JC Vaught \\
    \textit{Department of Mechanical Engineering} \\
    University of South Carolina
}
\date{\today}

% Document
\begin{document}

\maketitle

\section*{Introduction}
This document explores various natural systems and their mechanisms for regulating flow, receiving and storing energy and water, and transferring resources. These solutions from nature are presented alongside potential technological or mechanical analogs that could inspire sustainable engineering designs. 

\section*{Nature-Inspired Solutions for Each Function}
\begin{table}[h!]
    \centering
    \caption{Nature-Inspired Solutions for User Adjusting Flowrate}
    \begin{tabular}{p{4cm} p{5.5cm} p{5.5cm}}
        \toprule
        \textbf{Natural System} & \textbf{Description} & \textbf{Engineering Solution} \\
        \midrule
        Leaf Stomata Regulation & Leaves open and close stomata to regulate water vapor and gas exchange via guard cells. & Adjustable valves in irrigation systems that can be opened or closed manually. \\
        Plant Root Pressure Adjustment & Plants adjust root pressure to control water uptake from soil. & Manual pumps or pressure regulators for water pressure adjustment. \\
        Anemone Tentacle Movement & Sea anemones move tentacles to adjust water flow and capture prey. & Flexible tubes or nozzles that can be repositioned to direct flow. \\
        Frog Vocal Sac Inflation & Frogs control airflow to vocal sacs, regulating sound production. & Adjustable bladders or chambers to regulate flow. \\
        Earthworm Movement in Soil & Earthworms change soil structure as they move, influencing water infiltration. & Adjustable soil aerators or perforators to modify water flow into soil. \\
        \bottomrule
    \end{tabular}
    \label{tab:user_adjusting_flowrate}
\end{table}

\begin{table}[h!]
    \centering
    \caption{Nature-Inspired Solutions for Receiving Energy}
    \begin{tabular}{p{4cm} p{5.5cm} p{5.5cm}}
        \toprule
        \textbf{Natural System} & \textbf{Description} & \textbf{Engineering Solution} \\
        \midrule
        Plant Photosynthesis & Plants use chlorophyll to absorb sunlight, converting it into chemical energy. & Solar panels to capture solar energy and power system components. \\
        Wind Seed Dispersal & Seeds are dispersed by wind energy. & Wind turbines or small windmills generating electricity. \\
        Lightning Energy Transfer & Lightning transfers electrical energy from the atmosphere to the ground. & Static energy collectors or atmospheric electricity generators harnessing atmospheric energy. \\
        Frictional Heat from Movement & Tectonic plate movements generate heat through friction. & Mechanical devices converting kinetic energy into heat via friction. \\
        Magnetotactic Bacteria & Bacteria align with the Earth's magnetic field for orientation. & Magnetic field generators inducing current in coils to generate electricity. \\
        \bottomrule
    \end{tabular}
    \label{tab:receives_energy}
\end{table}

\begin{table}[h!]
    \centering
    \caption{Nature-Inspired Solutions for Storing Water}
    \begin{tabular}{p{4cm} p{5.5cm} p{5.5cm}}
        \toprule
        \textbf{Natural System} & \textbf{Description} & \textbf{Engineering Solution} \\
        \midrule
        Soil Absorption & Soil absorbs rainwater, storing it for plant use. & Infiltration basins or rain gardens to enhance groundwater recharge. \\
        Wetland Retention & Wetlands retain rainwater like natural sponges. & Artificial wetlands or retention ponds capturing runoff. \\
        Moss Water Retention & Mosses hold water in their tissues. & Absorbent materials like sponges in green roofs. \\
        Organic Matter Moisture Holding & Humus-rich soil retains moisture effectively. & Compost and mulch incorporated to enhance soil water retention. \\
        Fungal Mycelium Networks & Fungi distribute water through soil networks. & Mycorrhizal fungi to improve soil water storage. \\
        \bottomrule
    \end{tabular}
    \label{tab:stores_water}
\end{table}

\begin{table}[h!]
    \centering
    \caption{Nature-Inspired Solutions for Transferring Water}
    \begin{tabular}{p{4cm} p{5.5cm} p{5.5cm}}
        \toprule
        \textbf{Natural System} & \textbf{Description} & \textbf{Engineering Solution} \\
        \midrule
        Plant Transpiration & Plants release water vapor, adding moisture to the atmosphere. & Greenhouse systems where transpiration contributes to humidity, collected via condensation. \\
        Groundwater Flow & Water moves through underground aquifers. & Subsurface irrigation systems moving water through soil layers. \\
        Dew Deposition & Dew transfers atmospheric moisture to surfaces. & Surfaces optimized for dew collection, directing water to the system. \\
        Water-Based Seed Dispersal & Seeds are transported by water to new locations. & Hydro-seeding techniques distributing seeds via water flow. \\
        Erosion Water Movement & Water carries soil and nutrients across land. & Erosion control measures directing water flow strategically. \\
        \bottomrule
    \end{tabular}
    \label{tab:transfers_water}
\end{table}

\begin{table}[h!]
    \centering
    \caption{Nature-Inspired Solutions for Adjusting Flowrate}
    \begin{tabular}{p{4cm} p{5.5cm} p{5.5cm}}
        \toprule
        \textbf{Natural System} & \textbf{Description} & \textbf{Engineering Solution} \\
        \midrule
        Leaf Stomata Response & Stomata respond automatically to light and humidity changes. & Automatic valves controlled by sensors adjusting flowrate based on conditions. \\
        Beaver Dam Water Level Control & Beavers modify dams to regulate water levels. & Adjustable overflow outlets to maintain consistent water levels. \\
        Floodplain Water Absorption & Floodplains absorb excess water, reducing downstream flow. & Buffer zones or holding areas storing excess water temporarily. \\
        Vascular Nutrient Flow & Plants regulate nutrient flow through vascular tissues. & Nutrient dosing systems adjusting flow based on plant uptake. \\
        Plant Growth Hormones & Hormones regulate plant growth and sap flow. & Growth-based sensors adjusting irrigation based on growth stages. \\
        \bottomrule
    \end{tabular}
    \label{tab:adjust_flowrate}
\end{table}

\begin{table}[h!]
    \centering
    \caption{Nature-Inspired Solutions for Delivering Water}
    \begin{tabular}{p{4cm} p{5.5cm} p{5.5cm}}
        \toprule
        \textbf{Natural System} & \textbf{Description} & \textbf{Engineering Solution} \\
        \midrule
        Rainfall Distribution & Rain delivers water over wide areas. & Overhead sprinklers distributing water evenly across the garden. \\
        Animal Urine Distribution & Animals distribute water and nutrients through urination. & Fertigation systems delivering water and nutrients directly to soil. \\
        Floodplain Irrigation & Floods spread water and enrich soil. & Flood irrigation allowing water to flow over the soil surface. \\
        Wave Action & Waves carry water onto shores. & Oscillating sprinklers simulating wave motion. \\
        Nectar Secretion & Plants provide nectar, delivering water and sugars to pollinators. & Foliar sprays delivering water and nutrients directly to leaves. \\
        \bottomrule
    \end{tabular}
    \label{tab:deliver_water}
\end{table}

\end{document}
